% HD-S44. GPT-6 Astra Pro /root/sixfolds; inherited effort not exposed.
% 21 November 2026. Exact rank-three support and cubic-field descent.
\subsection{HD-S44: codimension-three supports and cubic blocks}
\label{sec:codimension-three-supports}

Let \(K\subset\mathbb C\) have characteristic zero, let
\(G\subset GL_K(V)\) be finite and absolutely irreducible,
and let \(1\ne N\triangleleft G\) be generated by actual
elements \(x\) with \(\operatorname{rank}(x-I)=3\).
Take all \(G\)-conjugates of a generating set, denoted
\(\mathcal X\). The following structural assertion is over
the specified field, before any imaginary-quadratic application.

The complex \(N\)-isotypic components have common multiplicity
one or three. At multiplicity one, either all components are
\(K\)-defined, or they form disjoint \(K\)-defined triples
with one cyclic cubic projector field \(E/K\). On a triple
the supported subgroup is an actual irreducible \(E\)-reflection
group. At multiplicity three, all components are \(K\)-defined
and admit actual coordinates \(W_K\otimes_K K^3\), with the
supported subgroup an irreducible \(K\)-reflection group on
\(W_K\). Every asserted decomposition and field action is
constructed below.

In the imaginary-quadratic sixfold setting, the cubic-triple
case with reflection rank \(r\ge5\) has period rank at most
three. For \(r\le4\) it has an actual coefficient model of
rank \(3r\le12\), with unchanged ordinary Hermitian determinant.
Thus over the Gaussian field and the seven fields
\(\mathbb Q(\sqrt{-d})\), \(d=2,7,11,19,43,67,163\),
this case gives only the determinant classes already proved
in HD-S40. The individually descended and multiplicity-three
cases are structural reductions here; their further endgames
are not asserted.

\paragraph{1. Weighted supports and the cubic projector field.}
Finite averaging gives invariant complements. Write
\[
 V_{\mathbb C}=\bigoplus_{i\in I}U_i,\qquad
 U_i=W_i\otimes\mathbb C^e,\qquad \dim W_i=r.
 \tag{S44.1}
\]
Normality and absolute irreducibility make \(G\) transitive
on the irreducible \(N\)-types, so the dimensions and
multiplicities are common. For their characters form
\[
 p_i=\frac r{|N|}\sum_{n\in N}\chi_i(n^{-1})\rho(n).
 \tag{S44.2}
\]
Character and matrix-coefficient averaging show that these
are the complete primitive central projectors of the complex
image algebra. Their entries lie in a finite cyclotomic
Galois extension \(L/K\). Galois permutes them and commutes
with the permutation of \(G\), since the original matrices
belong to \(M_m(K)\).

For \(x\in\mathcal X\), put
\(a_i(x)=\operatorname{rank}((x-I)|_{W_i})\). Then
\[
                3=e\sum_i a_i(x).                         \tag{S44.3}
\]
It follows that \(e=1\) or \(3\). The weighted rank
function is individually invariant under Galois: apply an
automorphism to every minor of the literal matrix
\((x-I)p_i\), whose rank is \(e\,a_i(x)\).

Suppose first that \(e=1\), and let \(\Gamma\) be the
image of Galois on \(I\). A permutation commuting with a
transitive group and fixing one label fixes every label.
Thus \(\Gamma\) acts freely on \(I\). Each of its orbits
has size \(|\Gamma|\), and each rank function \(a_i(x)\)
is constant on an orbit. Equation (S44.3) therefore gives
\[
                            |\Gamma|\mid3.                 \tag{S44.4}
\]
If \(\Gamma=1\), all \(p_i\) have entries in \(K\).
Otherwise \(\Gamma=C_3\), and each generator support is
exactly one free orbit of three labels, with rank one on
each. Every label occurs in a support: transitivity would
otherwise make \(N\) trivial. Hence these supports are a
perfect system of triples, with possible repetitions.

Let \(E\) be the fixed field of the kernel of the Galois
permutation action. Then \(E/K\) is cyclic of degree three.
Write \(\tau\) for a generator, choose a projector \(p\),
and put \(p_j=\tau^j(p)\) and \(U=(p_0+p_1+p_2)V_K\).
The actual field action and comparison maps are
\[
 \begin{aligned}
 \iota:E&\longrightarrow\operatorname{End}_K(U),
 &\iota(a)&=\sum_{j=0}^2\tau^j(a)p_j,\\
 U_E&\longrightarrow pE^m,
 &u&\longmapsto pu,\qquad
 w\longmapsto w+\tau(w)+\tau^2(w).
 \end{aligned}                                             \tag{S44.5}
\]
Here \(U_E\) means \(U\) with the displayed \(E\)-action.
The first matrix is fixed by Galois, is multiplicative and
injective, and sends one to \(I_U\). The second two maps
are inverse by the orthogonal idempotent equations. The
inverse takes values in \(K^m\). In particular \(\dim_EU=r\).

Generators on distinct triples have disjoint supports and
commute. Their supported subgroups have trivial intersections,
so \(N\) is their actual direct product. The factor on
\(U\), denoted \(B\), acts faithfully on the \(E\)-space
in (S44.5): identity on the first component implies identity
on all three by \(\tau\). Its specified generators are
reflections, and it is absolutely irreducible there. For its
actual character field \(F=\mathbb Q(\operatorname{tr}_E B)\),
\[
                               KF=E.                       \tag{S44.6}
\]
Indeed \(KF\) is an intermediate field of the prime-degree
extension. If \(KF=K\), its three conjugate characters
would coincide, contrary to the three distinct \(N\)-types.
The triple stabilizer \(H\) in \(G\) normalizes \(B\)
and acts \(E/K\)-semilinearly. Its image in
\(\operatorname{Gal}(E/K)\) is all of \(C_3\), because
transitivity on the constituent labels supplies an element
carrying \(p\) to \(\tau(p)\).
The action of \(H\) on \(U\) is absolutely irreducible:
a proper invariant complex subspace, translated over the
triple blocks, would give a proper \(G\)-invariant subspace.

\paragraph{2. Multiplicity three and its literal tensor map.}
If \(e=3\), each generator has one support and acts on the
corresponding \(W_i\) as a reflection. For a supported
generator \(x\), the actual subspace
\[
                 U_i(K)=K[N]\operatorname{im}_K(x-I)        \tag{S44.7}
\]
has complexification \(W_i\otimes\mathbb C^3\):
the nonzero reflection line spans \(W_i\) under \(N\),
while the whole multiplicity factor is already in the image.
These are therefore individually \(K\)-defined components.
Their supported subgroups again form a direct product.

On one component let \(A\) be the \(K\)-span of its
supported subgroup. Matrix-coefficient averaging gives
\(A_{\mathbb C}=M_r(\mathbb C)\otimes I_3\). Thus \(A\)
is central simple. Its irreducible character is one third
of the trace on \(U\), hence lies in \(K\). For a
supported reflection its nontrivial eigenvalue is
\(\lambda=\chi(x)-r+1\in K\). Consequently
\[
 P=\frac{x-I_U}{\lambda-1}\in A,\qquad P^2=P,\qquad
 PAP=KP,\qquad \dim_K PU=3.
 \tag{S44.8}
\]
The identities follow after scalar extension from the
rank-one reflection projector, and hence hold over \(K\).
Put \(W_K=AP\) and \(M_K=PU\). The left action identifies
\(A=\operatorname{End}_K(W_K)\): its kernel is zero because
\(APA=A\), and both dimensions are \(r^2\).
Choose \(a_j,b_j\in A\) with \(\sum a_jPb_j=1\).
The two actual inverse maps are
\[
 \begin{aligned}
 AP\otimes_K PU&\longrightarrow U,&aP\otimes u&\longmapsto au,\\
 U&\longrightarrow AP\otimes_K PU,&
 u&\longmapsto\sum_j a_jP\otimes Pb_j u.
 \end{aligned}                                             \tag{S44.9}
\]
The equation \(PAP=KP\) verifies both composites directly.
Thus the supported subgroup is an actual irreducible
\(K\)-reflection group on \(W_K\), acting identically on
the three-dimensional \(M_K\). As in (S42.11), transformed
matrix units give an explicit \(K\)-inner implementation
of each algebra automorphism; removing it leaves an operator
on \(M_K\). The component stabilizer therefore has actual
projective factor homomorphisms to \(PGL(W_K)\) and
\(PGL(M_K)\). This concludes the structural assertion.

\paragraph{3. The cubic case's actual rank-three axes.}
We now impose the imaginary-quadratic exact-character-field
sixfold conventions and return to the cubic case.
For its prescribed positive invariant coefficient form,
character unitarity and \(n\mapsto n^{-1}\) in (S44.2)
give \(p_i^\dagger=p_i\). Thus the triple sums are its
actual orthogonal projectors.
For \(r\ge5\), Cohen's published classification leaves the
real rational types, the imprimitive \(G(\ell,e,r)\), and
the two nonreal primitive groups \(G_{33},G_{34}\).
The precise locators are Theorem(2.4), pp.386--387;
(4.4)(i), p.399; table(4.16), p.412; and Theorem(5.12),
p.435, in \emph{Finite complex reflection groups},
Ann.\ Sci.\ ENS9(1976),379--436.
The primary article is
\url{https://www.numdam.org/item/ASENS_1976_4_9_3_379_0.pdf}.
The real rank-at-least-five types have rational root models.
For \(G_{33},G_{34}\), the actual roots on p.409 have a
common factor \(1/\sqrt2\) and otherwise coordinates in
\(\mathbb Q(\omega)\). The reflection formula
\(s_v=I-2v\bar v^t/(\bar v^tv)\) therefore gives matrices
in \(\mathbb Q(\omega)\). The rank-five invariant subspace
is \((\epsilon_5-\epsilon_6)^\perp\), with a rational basis.
Thus their character fields are contained in
\(\mathbb Q(\omega)\). These alternatives give
\([KF:K]\le2\), contradicting (S44.6). This verifies the
field exclusion from the actual published matrices.

Hence \(B\) is imprimitive. For its standard model
\(G(\ell,e,r)\), let \(\zeta=\zeta_\ell\).
The actual diagonal elements
\(\operatorname{diag}(\zeta,\zeta^{-1},1,\ldots)\) and
\(\operatorname{diag}(\zeta,\zeta,\zeta^{-2},1,\ldots)\)
have traces giving
\[
 c=\zeta+\zeta^{-1},\quad v=2\zeta+\zeta^{-2},\qquad
 \zeta=\frac{v-c^2+1}{2-c}.
 \tag{S44.10}
\]
Here \(\ell=1\) is the reducible permutation representation
and does not occur. Thus \(c\ne2\), and these equations
prove \(F=\mathbb Q(\zeta)\), \(E=K(\zeta)\). The
complex equivalence is replaced by the actual \(E\)-linear
intertwiner: solve its finite linear equations over \(E\).
Their complex solution space has dimension one and contains
an invertible map, so every nonzero \(E\)-solution is
invertible. Every standard test element in (S44.10) is
pulled back by that exact intertwiner.

The diagonal subgroup \(D\) is the unique largest abelian
normal subgroup of \(B\), because \(B/D=S_r\) has no
nontrivial abelian normal subgroup for \(r\ge5\).
Therefore every semilinear \(H\)-normalizer preserves \(D\).
Its distinct coordinate characters \(\lambda_j:D\to E^*\)
give actual \(E\)-linear projectors
\[
 f_j=|D|^{-1}\sum_{d\in D}\lambda_j(d)^{-1}\rho_E(d).
 \tag{S44.11}
\]
Applying the field action (S44.5) to the coefficients
makes these literal \(K\)-linear projectors on \(U\);
their images have \(K\)-dimension three and are permuted
transitively by \(H\), and consequently by \(G\) over
all original triples.

These images are orthogonal for the prescribed coefficient
form. Given distinct axes \(j,k\), choose a third axis \(t\).
The actual diagonal element with \(\zeta\) at \(j\),
\(\zeta^{-1}\) at \(t\), and one elsewhere is the identity
on axis \(k\). Its difference from the identity is
invertible on the \(K\)-space of axis \(j\), since all
three conjugates of \(\zeta\ne1\) remain nonidentity.
Invariance of the form then makes the two axes orthogonal.

Let \(J\) stabilize one axis, and write \(\sigma\) for
its faithful image there. This representation is absolutely
irreducible: \(D\) has three pairwise distinct characters
on its complexification, since the values of \(\lambda_j\)
include \(\zeta\) and \(K(\zeta)=E\). An element of \(H\)
inducing \(\tau\), composed with an internal coordinate
permutation, stabilizes this axis and cycles those three
characters. The only invariant sums of their eigenspaces
are zero and the whole space. The original component
translates therefore give the actual induced representation
\(\operatorname{Ind}_J^G\sigma=V\); in particular its
exact character field is \(K\).

On this one-dimensional \(E\)-axis, every element is literally
\[
                         u\longmapsto a\,\tau^j(u),
                         \qquad a\in E^*,\quad j=0,1,2.
 \tag{S44.12}
\]
For \(j=0\), a nonidentity element has zero fixed space.
For \(j=1,2\), two nonzero fixed vectors have ratio fixed
by \(\tau^j\), hence ratio in \(K\). Its fixed space
therefore has \(K\)-dimension at most one. Every nonidentity
element of this faithful rank-three image has codimension
at least two. In particular it has no reflection.

\paragraph{4. Quotients, original corners, and the period bound.}
Here are the actual polarized maps used both for the cubic
triple and for its descended axes. For any such \(K\)-block
\(T\), let \(J\) be its stabilizer, \(\sigma\) its action,
and \(J_0=\ker\sigma\). For the original connected cover put
\[
 Y=\widetilde C/J,\qquad C'=\widetilde C/J_0,\qquad
 \alpha:\widetilde C\longrightarrow C'.
 \tag{S44.13}
\]
The induction map is \(t\otimes u\mapsto\rho(t)u\).
The corresponding pushforward functions are
\(F_v(x)=\pi_T\rho(x)v\), with \(F(hx)=\sigma(h)F(x)\);
their inverse is \(\sum_{J\backslash G}\rho(x)^{-1}F(x)\).
Right translation gives the original monodromy. The inversion
label \(Jt^{-1}\leftrightarrow tJ\) fixes the coset convention.
These maps commute with taking each local inertia's invariant
subspace and hence extend across every branch stalk.

For a coefficient vector in \(T\), its normalized rational
matrix-coefficient idempotents satisfy \(e_G=e_J\):
off-\(J\) coefficients vanish by the orthogonal block
decomposition and \(\dim V/|G|=\dim T/|J|\).
Pullback of the associated primitive Jacobian factors is
therefore the actual isogeny and obeys
\[
 \alpha_*\alpha^*=[|J_0|],\qquad
 (\alpha^*)^*\theta_{J\widetilde C}=|J_0|\theta_{JC'}.
 \tag{S44.14}
\]
For a prescribed original primitive line, use exactly the
integer-cleared matrix-coefficient map \(M b_{v,w}\) of
(S41.15)--(S41.16), with multiplier
\(M^2h(v,v)/h(w,w)\). Thus the original polarized period
map is retained up to an actual constant isogeny, and the
ordinary rank-six Hermitian determinant is unchanged:
every positive rational multiplier has sixth power a norm.
Induction of characters, together with the matrices over
\(K\), proves the exact coefficient field assertion above.

For the rank-three axis the parabolic dimension is six,
and the exact Euler formula is
\[
              6=3(2g(Y)-2)+\sum_a c_a,\qquad c_a\ge2.
 \tag{S44.15}
\]
If \(g(Y)=0\), there are at most six true coefficient
branch values and the moduli dimension is at most three.
If \(g(Y)=1\), there are at most three such values and
the pointed moduli dimension is at most three; the
unpointed case is excluded by (S44.15). If \(g(Y)=2\),
there is no coefficient branch and the base moduli dimension
is three. Higher genus is impossible. For fixed finite
monodromy data, the covers over a marked base form a finite
discrete set; passing to a finite local parameter cover
therefore gives an actual family map from these pointed-curve
moduli. The polarized period rank is at most three in every
case. The exact quotient and corner maps above give the same
bound for the original sixfold.

For \(r\le4\), apply (S44.13)--(S44.14) to the whole
cubic triple instead. Its faithful absolutely irreducible
coefficient image has rank \(3r\le12\) and exact field
\(K\). At period rank nine its base must have genus zero:
genus one has at most six parameters, genus two permits
only rank three and has dimension three, and higher genus
is impossible. Thus the already proved actual coefficient
rank-at-most-fourteen result HD-S40 applies over its stated
eight fields. This proves the claimed cubic-case conclusion.
